\paragraph{Reflection by witness closure and construction of the image.}
The Collection--Replacement step of Lemma~\ref{thm:ch-6b} admits the following
explicit proof. We retain the hierarchy already generated from the operator
code and the realized record, and abbreviate
\(\mathfrak G=\mathfrak G_{\rm HMT}(h,m_\infty)\).
The argument uses Separation and Replacement in the ambient metatheory,
without presupposing these schemes in \(\mathfrak G\).

Let \(\Phi\) be a finite collection of formulas. First translate their
universal quantifiers using negation and existential quantification, and
close the resulting collection under subformulas. For every initial level
there is a later limit ordinal \(\theta\) such that
\[
 \mathfrak G_\theta\models\varphi(\bar a)
 \quad\Longleftrightarrow\quad
 \mathfrak G\models\varphi(\bar a)
 \qquad
 (\varphi\in\Phi,\ \bar a\in\mathfrak G_\theta^{<\omega}).
\]
For each fixed formula, satisfaction in the class is interpreted by
relativizing its quantifiers to the definable class \(\mathfrak G\);
no uniform truth predicate is introduced.

\begin{proof}
Fix \(\alpha\). For each existential subformula
\(\exists y\,\psi(y,\bar x)\) of \(\Phi\), the tuples
\(\bar a\in\mathfrak G_\alpha^{<\omega}\) of the corresponding arity
that have a witness in \(\mathfrak G\) form a set by Separation in the
metatheory. Assign to each tuple the least ordinal \(\beta\) for which
some \(y\in\mathfrak G_\beta\) satisfies
\(\mathfrak G\models\psi(y,\bar a)\). This ordinal exists because
\(\mathfrak G\) is the union of its levels. Replacement in the metatheory
bounds these first stages of witnesses. Since \(\Phi\) is finite, a
single bound \(f(\alpha)>\alpha\), valid for all its existential
subformulas, can be chosen definably. Thus every tuple under consideration
has \emph{at least one witness} in \(\mathfrak G_{f(\alpha)}\).

Take \(\alpha_0\) large enough to contain the initial parameters and set
\[
 \alpha_{n+1}=f(\alpha_n),\qquad
 \theta=\sup_{n<\omega}\alpha_n.
\]
By continuity of the hierarchy, every finite tuple from
\(\mathfrak G_\theta\) belongs to some \(\mathfrak G_{\alpha_n}\).
If an existential formula of \(\Phi\) with that tuple is true in
\(\mathfrak G\), it has a witness in
\(\mathfrak G_{\alpha_{n+1}}\subseteq\mathfrak G_\theta\).
Induction on formulas now proves the equivalence: atomic formulas are
preserved in the induced structure; connectives preserve the equivalence;
the existential step uses that witness in one direction and the same
witness, viewed in \(\mathfrak G\), in the other. The prior translation
of universal quantifiers also includes the existential formulas expressing
their possible counterexamples. This proves the required finite reflection.
\end{proof}

Now let \(F(x,y,\bar p)\) be functional on \(a\) in \(\mathfrak G\).
Apply the construction to the formulas expressing \(F\), the existence
of its witnesses, and its image, with \(a,\bar p\in\mathfrak G_\theta\).
Transitivity gives \(a\subseteq\mathfrak G_\theta\). Every
\(x\in a\) has its value within that level, and reflection identifies
the exact image as
\[
 b=\{y\in\mathfrak G_\theta:
       \mathfrak G_\theta\models
       \exists x\in a\ F(x,y,\bar p)\}.
\]
This set is definable over \(\mathfrak G_\theta\), with parameters from
the same level. By the successor operation of the hierarchy,
\(b\in\operatorname{Def}(\mathfrak G_\theta)
\subseteq\mathfrak G_{\theta+1}\). This establishes Replacement in
\(\mathfrak G\). If the existence of witnesses is retained and uniqueness
is omitted, the same construction provides a set of witnesses for Collection.

\paragraph{Arity of names and internal selection.}
In the well-order of Lemma~\ref{thm:ch-6b}, the enumeration of the base
assigns to each element its least natural-number code. At each successor
stage, formula codes are ordered first; once a code is fixed, the arity
of its parameter tuple is fixed. Lexicographic order is then applied to a
\emph{finite product of fixed arity} of the previously constructed
well-orders. The least name represents each new value, and limit levels
retain the order of first appearance. This specifies the ordering of
tuples without combining different arities into a single lexicographic
order. Selecting the least element of every nonempty set, together with
Replacement already established, forms internally the choice functions
used subsequently. This construction of names precedes the cardinal
injections and receives no bijection of the continuum as data.
